\documentclass{report}
\usepackage{amsmath,amssymb}
\setlength{\parindent}{0mm}
\setlength{\parskip}{1em}
\begin{document}
\begin{center}
\rule{6in}{1pt} \
{\large Luis Garcia Ramos \\
%-orig: {\bf On the spectrum of deflated matrices with applications to the deflated shifted Laplace preconditioner for the Helmholtz equation}
{\bf The spectrum of deflated matrices and the deflated shifted laplace preconditioner for the Helmholtz equation}
}

Technische Universit\"{a}t Berlin \\ Institut f\"ur Mathematik \\ Stra\ss e des 17 Juni 136 \\ 10623 Berlin \\ Germany
\\
{\tt garcia@math.tu-berlin.de}\\
Reinhard Nabben\end{center}

The deflation technique to accelerate Krylov subspace iterative methods
for the solution of linear systems has been known for a long time. The
first landmark papers are due
to Nicolaides and Dost\'{a}l in the late eighties where deflation is used
for Hermitian positive definite (HPD) linear systems. In the last decade
deflation was used and analyzed in combination with domain decomposition
and multigrid methods, resulting in very effective algorithms. Examples
are the multilevel Krylov methods introduced in [1,2], see also [3,4]
where multilevel deflation techniques are presented. Although these
algorithms work very well in practice for non-Hermitian problems, not
many theoretical results are known so far in this direction. Here we show
inclusion regions for the spectrum of an arbitrary deflated matrix based
on the field of values, which generalize known results for HPD systems.
Moreover, for deflated GMRES we show a residual bound based on the field
of values. We apply our results to linear systems arising from the
Helmholtz equation. We focus on the combination of the complex shifted
Laplace (CSL) preconditioner [5] with the multilevel Krylov technique.
Numerical examples indicate that the eigenvalues of the deflated
CSL-preconditioned system lie on exact the same circles as the
CSL-preconditioned linear systems and are shifted away from zero. Here we
are able to prove these surprising results for any wavenumber and any
dimension using M\"{o}bius transformations and the Spectral Mapping
Theorem. Our new results help to explain the good performance of
multilevel Krylov methods for the Helmholtz equation.

References

[1] Y. A. Erlangga and R. Nabben, Multilevel projection-based nested
Krylov iteration for boundary value problems, SIAM Journal on Scientific
Computing, 30 (2008), pp. 1572�1595.

[2] Y. A. Erlangga and R. Nabben, On a multilevel Krylov method for the
Helmholtz equation preconditioned by shifted Laplacian, Electronic
Transactions on Numerical Analysis, 31 (2008), pp. 403�424.

[3] A. H. Sheikh, D. Lahaye, L. Garc\'{i}a Ramos, R. Nabben, and C. Vuik,
Accelerating the Shifted Laplace Preconditioner for the Helmholtz
Equation by Multilevel Deflation, (2015), preprint.

[4] A. H. Sheikh, D. Lahaye, and C. Vuik, On the convergence of shifted
Laplace preconditioner combined with multilevel deflation, Numerical
Linear Algebra with Applications, 20 (2013), pp. 645�662.

[5] Y. A. Erlangga, C. Vuik, and C. W. Oosterlee, A novel multigrid-based
preconditioner for heterogeneous Helmholtz problems, SIAM Journal on
Scientific Computing, 27 (2006), pp. 1471�1492.


\end{document}
